\section*{Part XI: The Founder's Governance Operating System}
\addcontentsline{toc}{section}{Part XI: The Founder's Governance Operating System}

\subsection{69. The weekly review}

Review cash, looming statutory deadlines, promises made in sales, payment and delivery risks, privacy/security incidents and unresolved conflicts. Escalate urgent compliance and safeguarding matters immediately rather than waiting for a meeting.

\subsection{70. The monthly review}

Reconcile accounts; review revenue quality, tax liabilities, contractual commitments, employment matters, customer outcomes, vendor balances, receivable aging, material incidents, and decision logs. Management should explain deviations and corrective actions with named owners.

\subsection{71. Quarterly board discipline}

Issue a concise packet in advance. Record material approvals, conflicts, significant risks, capital allocation and legal changes. Reserve a session to ask what the management team might be systematically misunderstanding. If there is no board beyond statutory minimum, create a founder review with an external financially literate adviser for high-consequence decisions.

\subsection{72. Annual governance health check}

Check statutory filings, audit quality, shareholder records, board and committee functioning, contracts, related-party transactions, insurance, cybersecurity, academic and research standards, workforce processes, grant reporting, succession and financial reserves. Compare the company's risk and organizational complexity to its controls.

\subsection{73. The governance stop rule}

Do not sign an equity or debt agreement whose rights you cannot model; acquire property with unresolved title; claim degree authority without legal basis; take restricted donations without compliance; publish scientific claims without evidence; or use learner/employer data without proper rights. Escalate to qualified advisers and postpone the transaction when these gates are unresolved. A delayed deal is often cheaper than a disputed 1.
